\chapter{光谱与光的信息}\label{ch:g10:light-spectra}

一束星光是我们永远能够掌握的唯一恒星样本，而且它
竟是一份相当慷慨的样本。把这束光展开成光谱——
玻璃棱镜可以做到，一帘雨滴也可以做到——光便
变成一条信息：连续的彩虹色背景揭示恒星的
温度，而细线组成的“条码”逐一拼出恒星
由哪些元素构成。本章将学习读懂这信息的两部分——
先看实验室里的灯，再看太阳。

\section{白光及其光谱}

\begin{definition}[白光的分解]\label{def:g10:light-spectra:white}
当一束细窄的太阳光穿过玻璃棱镜时，射出后的光会散开
成一条由红到紫的连续色带。这条色带就是
光的\emph{光谱}\index{光谱}。包含全部可见颜色的
光（如太阳光或普通灯泡的光）称为
\emph{白光}\index{白光}。只含单一纯色的光，
即棱镜无法再分解的光，称为
\emph{单色光}\index{单色光}（例如激光束）；
由几种颜色混合而成的光，称为
\emph{复色光}\index{复色光}。能完成这一分解并读出光的组成的仪器，
称为
\emph{分光镜}\index{分光镜}。
\end{definition}

\begin{omfigure}
\begin{tikzpicture}[scale=1.0, font=\small]
  \draw[thick, fill=blue!6] (0,0) -- (1.4,2.6) -- (2.8,0) -- cycle;
  \draw[line width=1.6pt] (-2.4,1.95) -- (0.75,1.4)
    node[pos=0.22, above] {白光};
  \draw[gray] (0.75,1.4) -- (2.05,1.05);
  \draw[red, thick]              (2.05,1.05) -- (5.4,0.30);
  \draw[orange, thick]           (2.05,1.05) -- (5.36,0.08);
  \draw[yellow!75!black, thick]  (2.05,1.05) -- (5.31,-0.14);
  \draw[green!65!black, thick]   (2.05,1.05) -- (5.25,-0.36);
  \draw[blue, thick]             (2.05,1.05) -- (5.17,-0.58);
  \draw[violet, thick]           (2.05,1.05) -- (5.08,-0.80);
  \node[right] at (5.4,0.30)  {红色（偏折最小）};
  \node[right] at (5.08,-0.80) {紫色（偏折最大）};
\end{tikzpicture}

{\small 棱镜将一束白光扇形展开为光谱；在
屏幕上，这束光扇形映出一条连续的彩虹色带。}
\end{omfigure}

\begin{remark}[棱镜与雨滴]\label{rem:g10:light-spectra:rainbow}
棱镜之所以起作用，是因为玻璃对紫光的偏折比对红光略大
——\emph{为何}如此，属于折射的研究内容，见
\cref{ch:g10:refraction}；这里只使用这一效应。水滴也起着同样的作用，
即经过的阵雨中每一滴雨水都相当于一个小棱镜，接收
白色阳光，并按颜色将其分开返回。彩虹便是
太阳的光谱，铺展在天空中。
\end{remark}

\begin{definition}[波长：辐射的“身份证”]\label{def:g10:light-spectra:wavelength}
每一种单色辐射都由一个数来刻画，即其
\emph{波长}\index{波长} $\lambda$，以纳米为单位测量
（$\qty{1}{nm} = \qty{e-9}{m}$）。波长正是该辐射的“身份证”：
只要给出 $\lambda$，便说尽了
关于这一纯色的全部内容。人眼能响应的波长约在
$\qty{400}{nm}$（紫）至 $\qty{800}{nm}$（红）之间——\emph{可见
光谱}\index{可见光谱}。超出两端，辐射仍然存在，只是我们看不见：
\emph{紫外线}\index{紫外线}位于 $\qty{400}{nm}$ 以下，
\emph{红外线}\index{红外线}位于 $\qty{800}{nm}$ 以上。
波长究竟度量什么——光波上一个波纹的长度——
这一点将在\cref{ch:g10:signals-and-waves}中说明；
本章只需关注这个数本身。
\end{definition}

\begin{example}[读取波长]\label{ex:g10:light-spectra:classify}
一支绿色激光笔发出 $\qty{532}{nm}$ 的光，属于可见光范围，
呈绿色。一只杀菌用汞灯，其发光波长为
$\qty{254}{nm}$：属于紫外线，人眼不可见（正因不可见，
它对眼睛才格外危险）。电视遥控器的二极管波长为 $\qty{940}{nm}$：
属于红外线——对准手机摄像头，其感光元件能看到略超过
$\qty{800}{nm}$ 的光，于是屏幕上便会出现这些不可见的闪光。
\end{example}

\section{连续光谱与温度}

\begin{definition}[连续光谱]\label{def:g10:light-spectra:continuous}
致密的炽热物质——灯泡的灯丝、熔融的钢水、恒星表面——会发光，
其光包含较宽范围内的\emph{所有}波长，
形成一条不间断、无任何缺失的色带。这样的
光谱称为\emph{连续光谱}\index{连续光谱}。它不携带
物体化学性质的任何痕迹：白热的铁棒与白热的铜棒
发光完全相同。它所编码的——完全而唯一地——
就是温度。
\end{definition}

\begin{definition}[热力学温度]\label{def:g10:light-spectra:kelvin}
对于辐射定律，\emph{热力学温度}\index{热力学
温度}不从水的凝固点起算，而是从
可能达到的最低温度状态起算，
即 $\qty{-273}{\degreeCelsius}$。单位是\emph{开尔文}\index{开尔文}
（符号 $\unit{K}$），两者的温度刻度大小相同：
\[
T \text{ (单位： } \unit{K}) = \theta \text{ (单位： } \unit{\degreeCelsius}) + 273 .
\]
因此，温度为 $\qty{20}{\degreeCelsius}$ 的房间，热力学温度为 $\qty{293}{K}$，
而太阳表面温度约为 $\qty{5500}{\degreeCelsius}$，其热力学温度
接近 $\qty{5800}{K}$。
\end{definition}

\begin{proposition}[温度越高，越亮且越偏蓝（维恩定律）]\label{prop:g10:light-spectra:wien}
当致密发光物体的温度升高时，
\begin{enumerate}
  \item 它在\emph{每一}波长上的辐射都增强：整个光谱
        整体变亮；
  \item 辐射最强的波长 $\lambda_{\max}$
        向短波长（蓝色）一侧移动，
        满足
        \[
        \lambda_{\max} \times T = \qty{2.90e-3}{m.K},
        \]
        其中 $T$ 为热力学温度，单位为开尔文。
\end{enumerate}
\end{proposition}
\admitted

\begin{remark}[这一定律的来源]\label{rem:g10:light-spectra:wienorigin}
维恩定律概括了十九世纪末人们对熔炉和灯丝
所做的细致测量。推导它需要辐射的量子理论——
这一光谱恰恰是经典物理学\emph{无法}解释的，
也正是迫使普朗克提出量子的难题。
严格的推导在本科三年级
分册中完成。
\end{remark}

\begin{omfigure}
\begin{tikzpicture}[scale=1.0, font=\small]
  \foreach \y in {0, 1.1, 2.2}
    \foreach \xa/\xb/\c in {0/0.6/violet, 0.6/1.35/blue,
        1.35/2.4/green!75!black, 2.4/2.85/yellow, 2.85/3.45/orange,
        3.45/6/red}
      \fill[\c] (\xa,\y) rectangle (\xb,\y+0.55);
  \foreach \xa/\xb/\o in {0/0.6/0.85, 0.6/1.35/0.7, 1.35/2.4/0.5,
      2.4/2.85/0.35, 2.85/3.45/0.25, 3.45/6/0.1}
    \fill[black, opacity=\o] (\xa,0) rectangle (\xb,0.55);
  \foreach \xa/\xb/\o in {0/0.6/0.5, 0.6/1.35/0.35, 1.35/2.4/0.2,
      2.4/2.85/0.12, 2.85/3.45/0.08, 3.45/6/0.08}
    \fill[black, opacity=\o] (\xa,1.1) rectangle (\xb,1.65);
  \node[left] at (-0.15, 0.275) {$\qty{3000}{K}$};
  \node[left] at (-0.15, 1.375) {$\qty{4500}{K}$};
  \node[left] at (-0.15, 2.475) {$\qty{6000}{K}$};
  \draw[-Stealth] (0,-0.5) -- (6.7,-0.5)
    node[right] {$\lambda$ (\unit{nm})};
  \foreach \x/\l in {0/400, 1.5/500, 3/600, 4.5/700, 6/800}
    \draw (\x,-0.42) -- (\x,-0.58) node[below] {\l};
\end{tikzpicture}

{\small 三种温度下连续光谱的可见光波段：
光源越冷，波段越暗——蓝端
最先变淡。}
\end{omfigure}

\begin{example}[测定太阳的温度]\label{ex:g10:light-spectra:suntemp}
太阳的连续光谱在
$\lambda_{\max} = \qty{500}{nm}$ 附近强度最大。由维恩定律可求得其表面
温度：
\[
T = \frac{\qty{2.90e-3}{m.K}}{\qty{5.00e-7}{m}} = \qty{5800}{K},
\]
约 $\qty{5500}{\degreeCelsius}$——在相距 $150$ 百万千米之外，
借助棱镜和刻度测得。
\end{example}

\begin{omfigure}
\begin{tikzpicture}
\begin{axis}[omaxis, width=9.4cm, height=5.6cm,
    xlabel={$\lambda$ (\unit{nm})}, ylabel={强度（任意单位）},
    xmin=0, xmax=2150, ymin=0, ymax=9.5,
    xtick={400, 800, 1200, 1600, 2000}, ytick=\empty]
  \fill[gray, opacity=0.12] (axis cs:400,0) rectangle (axis cs:800,9.5);
  \node[gray, anchor=south] at (axis cs:600,8.4) {可见光};
  \addplot[omThm, thick, domain=180:2050, samples=160]
    {3e16 * x^(-5) / (exp(2400/x) - 1)};
  \addplot[omDef, thick, domain=180:2050, samples=160]
    {3e16 * x^(-5) / (exp(3600/x) - 1)};
  \draw[dashed, omThm] (axis cs:480,0) -- (axis cs:480,8.0);
  \draw[dashed, omDef] (axis cs:725,0) -- (axis cs:725,1.05);
  \node[omThm, anchor=south west] at (axis cs:620,6.4)
    {高温恒星，$\qty{6000}{K}$};
  \node[omDef, anchor=south west] at (axis cs:950,1.3)
    {低温恒星，$\qty{4000}{K}$};
\end{axis}
\end{tikzpicture}

{\small 两颗恒星的强度随波长的变化：较热的一颗
在每个波长上都辐射得更强，其峰值 $\lambda_{\max}$（虚线）
位于更短的波长处。}
\end{omfigure}

\section{受激气体的线状光谱}

\begin{definition}[发射光谱]\label{def:g10:light-spectra:emission}
低压气体在放电激发下（如霓虹灯管
或钠路灯）同样会发光，但经分光镜观察，其光
与彩虹全然不同。只有少数孤立的波长
出现，即黑色背景上的细亮线，称为\emph{谱
线}\index{谱线}。这样的光谱是一种\emph{发射
光谱}\index{发射光谱}\index{线状光谱}。
\end{definition}

\begin{example}[氢与钠]\label{ex:g10:light-spectra:hydrogen}
受激氢在可见光区恰好发出四条谱线：红线位于
$\qty{656}{nm}$，蓝绿线位于 $\qty{486}{nm}$，两条紫线分别位于
$\qty{434}{nm}$ 和 $\qty{410}{nm}$。受激钠发出的
基本是一条很强的黄橙色谱线，位于 $\qty{589}{nm}$——这正是
旧式路灯把整个街区都染成同一种颜色的原因：
其光中几乎别无其他成分。
\end{example}

\begin{proposition}[元素的“指纹”]\label{prop:g10:light-spectra:fingerprint}
每种化学元素都发出自己固定的一组谱线，
这组谱线在每个实验室、每个时代都相同，且没有两种元素
共享同一组谱线。因此，线状光谱能标识发光的元素，
正如指纹能标识一个人。
\end{proposition}
\admitted

\begin{remark}[为什么是谱线？]\label{rem:g10:light-spectra:whylines}
原子为何只发出某些特定波长，以及这些波长为何能
泄露元素的身份，要由原子的量子理论来解释，这在本分册的
最后一年中讲授：每种元素拥有一套能级阶梯，
每条谱线标记阶梯上一次跃迁。目前只需像侦探那样使用
这一“指纹”，无需追问手指是怎样长成的。
\end{remark}

\section{吸收光谱}

\begin{definition}[吸收光谱]\label{def:g10:light-spectra:absorption}
让白光\emph{穿过}一团低温低压气体，再将出射光色散：
连续彩虹重新出现，但被细黑
线刻划。气体已从透过的光中吸走了某些波长。
这样的光谱 --- 连续背景减去一组谱线 --- 称为
\emph{吸收光谱}\index{吸收光谱}。
\end{definition}

\begin{proposition}[发射谱线与吸收谱线重合]\label{prop:g10:light-spectra:kirchhoff}
气体恰好吸收它能够发射的波长：某元素吸收光谱中的暗
线，与其发射光谱中的亮线处于相同波长。
因此，
\cref{prop:g10:light-spectra:fingerprint}中的“指纹”既可在暗线中读出，
也可在亮线中读出。
\end{proposition}
\admitted

\begin{remark}[同一机制，两种表现]\label{rem:g10:light-spectra:mechanism}
这种巧合并非偶然：在 $\lambda$ 处吸收，所攀登的正是
在 $\lambda$ 处发射时所下降的同一能级阶梯 --- 本分册末尾的量子章节
将对此作精确阐述。该结论由基尔霍夫与本生
于 1859 年确立，是解锁下一节全部内容的钥匙。
\end{remark}

\begin{omfigure}
\begin{tikzpicture}[scale=1.0, font=\small]
  \fill[black] (0,1.5) rectangle (6,2.1);
  \draw[violet, line width=1.6pt] (0.15,1.5) -- (0.15,2.1);
  \draw[blue, line width=1.6pt]   (0.51,1.5) -- (0.51,2.1);
  \draw[cyan, line width=1.6pt]   (1.29,1.5) -- (1.29,2.1);
  \draw[red, line width=1.6pt]    (3.84,1.5) -- (3.84,2.1);
  \node[left] at (-0.15,1.8) {发射};
  \foreach \xa/\xb/\c in {0/0.6/violet, 0.6/1.35/blue,
      1.35/2.4/green!75!black, 2.4/2.85/yellow, 2.85/3.45/orange,
      3.45/6/red}
    \fill[\c] (\xa,0) rectangle (\xb,0.6);
  \foreach \x in {0.15, 0.51, 1.29, 3.84}
    \draw[black, line width=1.6pt] (\x,0) -- (\x,0.6);
  \node[left] at (-0.15,0.3) {吸收};
  \foreach \x in {0.15, 0.51, 1.29, 3.84}
    \draw[dashed, gray] (\x,0.6) -- (\x,1.5);
  \draw[-Stealth] (0,-0.5) -- (6.7,-0.5)
    node[right] {$\lambda$ (\unit{nm})};
  \foreach \x/\l in {0/400, 1.5/500, 3/600, 4.5/700, 6/800}
    \draw (\x,-0.42) -- (\x,-0.58) node[below] {\l};
\end{tikzpicture}

{\small 氢的两种光谱：亮的发射谱线（上）与暗的
吸收谱线（下）位于完全相同的波长处，即 $410$、$434$、
$486$ 与 $\qty{656}{nm}$。}
\end{omfigure}

\begin{example}[盛有钠蒸气的烧瓶]\label{ex:g10:light-spectra:sodiumflask}
白光穿过一只盛有低温钠蒸气的烧瓶。透射光谱为完整的彩虹，
只在 $\qty{589}{nm}$ 处有一条暗线 --- 正是钠灯最亮之处。
在暗室中从侧面观察烧瓶，可看到被“偷走”的光又被重新
发射出来：在同一波长处，一道微弱的
黄橙色光。
\end{example}

\section{光谱分析：读取恒星}

\begin{definition}[光谱分析]\label{def:g10:light-spectra:analysis}
\emph{光谱分析}\index{光谱分析}是指识别光源（火焰、灯、恒星）
中所含的化学元素，
其做法是将光谱中的谱线与实验室中
各元素的谱线目录逐一比对。
\end{definition}

\begin{method}[读取恒星光谱]\label{met:g10:light-spectra:reading}
恒星是一团由致密的热物质（表面）构成的球体，外面包着一层
更稀薄、较冷的大气层。因此，其光以连续光谱的形式到达，
其上布满吸收谱线 --- 且须对各部分分别读取：
\begin{enumerate}
  \item 找出连续背景的峰值 $\lambda_{\max}$；
        由维恩定律（\cref{prop:g10:light-spectra:wien}）可得表面
        温度 $T = \qty{2.90e-3}{m.K} / \lambda_{\max}$。
  \item 列出暗线的波长。
  \item 将其与目录相比对：某元素只有在其
        \emph{全部}强谱线均出现时，才可判定其存在。
  \item 得出结论：吸收发生在恒星的大气层中，故所识别的
        元素即为大气层中的元素。
\end{enumerate}
\end{method}

\begin{example}[太阳的“配方”]\label{ex:g10:light-spectra:sunrecipe}
太阳光谱中有数千条暗线，由夫琅禾费
于 1814 年测绘。其中最强的有 $410$、$434$、$486$ 与 $\qty{656}{nm}$ ---
即完整的氢“指纹” --- 另有位于
$\qty{589}{nm}$ 处的钠线，以及 $393$ 与 $\qty{397}{nm}$ 处的一对钙线。太阳大气
以氢为主，另含少量地球岩石中常见的元素。现代测量
对这一结论作了修正：
太阳按质量计约有四分之三为氢。
\end{example}

\begin{remark}[首先在天空中发现的元素]\label{rem:g10:light-spectra:helium}
1868 年日食期间，太阳大气的光谱中出现了一条亮线，
位于 $\qty{587.6}{nm}$ 处，与任何已知元素都不吻合。
由\cref{prop:g10:light-spectra:fingerprint}可知，未知的“指纹”意味着未知的元素，
于是宣布发现一种新元素，并以希腊语“太阳”\emph{helios}为其命名：
氦。该元素直至 1895 年才在地球上被分离出来 ---
它是唯一一种先在太空中被发现、后才在地球上找到的化学元素。
\end{remark}

\section{练习}

\begin{exercise}[$\star$]\label{exo:g10:light-spectra:1}
对每种辐射，判断它是紫外线、可见光还是红外线，
若为可见光，写出其颜色：$\qty{254}{nm}$（汞灯）、
$\qty{480}{nm}$、$\qty{589}{nm}$、$\qty{656}{nm}$、$\qty{1064}{nm}$
（切割激光器）。
\end{exercise}

\begin{exercise}[$\star$]\label{exo:g10:light-spectra:2}
一束细窄的阳光射到一块玻璃棱镜上。
\begin{enumerate}
  \item 描述放在棱镜后方屏幕上出现的现象。
  \item 光带的哪一端偏折得最厉害？
  \item 将阳光换成绿色激光笔的光束。
        屏幕上会出现什么，这说明
        激光具有怎样的特性？
\end{enumerate}
\end{exercise}

\begin{exercise}[$\star$]\label{exo:g10:light-spectra:3}
将每个光源与其所产生的光谱类型对应起来（连续、线状
发射或吸收光谱）：（a）白炽灯的灯丝；
（b）霓虹灯广告管；（c）铸造厂中的熔融铁水；（d）钠蒸气路灯；
（e）到达地球的太阳光。
\end{exercise}

\begin{exercise}[$\star$]\label{exo:g10:light-spectra:4}
关于彩虹：
\begin{enumerate}
  \item 什么起到了棱镜的作用？
  \item 只有当太阳在观察者身后、前方有雨时才能看到彩虹：
        每个条件分别提供了什么？
  \item 解释为什么彩虹能证明太阳光是白光。
\end{enumerate}
\end{exercise}

\begin{exercise}[$\star$]\label{exo:g10:light-spectra:5}
一位铁匠加热一根铁钉：它先发出暗红色的光，然后变成明亮的橙色，
再变成几乎白色。
\begin{enumerate}
  \item 将这三个阶段按温度从低到高排序。
  \item 此序列体现了\cref{prop:g10:light-spectra:wien}中的
        哪两个效应？
  \item 论证“白热比红热更热”这一说法。
\end{enumerate}
\end{exercise}

\begin{exercise}[$\star\star$]\label{exo:g10:light-spectra:6}
太阳的连续光谱在 $\lambda_{\max} = \qty{500}{nm}$ 处达到峰值。
\begin{enumerate}
  \item 求太阳表面的温度，单位为开尔文。
  \item 将其换算为摄氏度。
  \item 检验峰值位于可见光范围内，
        并说明其颜色。
\end{enumerate}
\end{exercise}

\begin{exercise}[$\star\star$]\label{exo:g10:light-spectra:7}
一支蜡烛的火焰约在 $\qty{1800}{K}$ 下发光；一只电热炉
调到“保温”挡，温度约为 $\qty{500}{K}$。
\begin{enumerate}
  \item 分别计算每个光源的 $\lambda_{\max}$，并给出其所属波段。
  \item 电热炉明显发热，手放在其上方能感受到辐射，
        但它看起来却是黑色的。请解释。
  \item 蜡烛的峰值同样在可见光范围之外，
        为什么我们仍能看见火焰？
\end{enumerate}
\end{exercise}

\begin{exercise}[$\star\star$]\label{exo:g10:light-spectra:8}
实验室资料给出，可见光范围内：氢
$\{410, 434, 486, 656\}\,\unit{nm}$；钠 $\{589\}\,\unit{nm}$；氦
$\{447, 502, 588, 668\}\,\unit{nm}$。现分析两个放电灯。
灯 A 显示位于 $434$、$486$ 和 $\qty{656}{nm}$ 处的谱线，以及一条
靠近可见光边缘的微弱紫色谱线；灯 B 只显示一条位于
$\qty{589}{nm}$ 处的强谱线。确定每个灯中的气体，并
仔细论证灯 B 中为何不含氦。
\end{exercise}

\begin{exercise}[$\star\star$]\label{exo:g10:light-spectra:9}
将白光射入一瓶低温钠蒸气中。
\begin{enumerate}
  \item 描述透射光的光谱。
  \item 其暗线位于哪个波长处，为什么恰好在
        这一波长处？
  \item 关闭灯并使房间变暗，蒸气现在
        由放电激发。描述新的光谱。
\end{enumerate}
\end{exercise}

\begin{exercise}[$\star\star$]\label{exo:g10:light-spectra:10}
太阳的高分辨率光谱中出现位于 $410$、$434$、
$486$、$517$、$518$ 和 $\qty{656}{nm}$ 处的暗线。参考数据：氢
$\{410, 434, 486, 656\}\,\unit{nm}$；镁 $\{517, 518\}\,\unit{nm}$。
\begin{enumerate}
  \item 这些谱线揭示了哪些元素？
  \item 这些谱线产生于太阳的哪一部分，
        以及它们为何呈暗线而非亮线？
\end{enumerate}
\end{exercise}

\begin{exercise}[$\star\star$]\label{exo:g10:light-spectra:11}
你的身体表面温度约为 $\qty{37}{\degreeCelsius}$。
\begin{enumerate}
  \item 换算为开尔文，并计算你自身热辐射的
        $\lambda_{\max}$。
  \item 它位于哪个波段？
  \item 推断为什么在黑暗房间里没有人会发出可见光，以及
        热像仪如何仍能找到人。
\end{enumerate}
\end{exercise}

\begin{exercise}[$\star\star\star$]\label{exo:g10:light-spectra:12}
白炽灯泡中的钨丝温度为 $\qty{2700}{K}$。
\begin{enumerate}
  \item 计算 $\lambda_{\max}$，并给出其所属的波段。
  \item 解释为什么这类灯泡作照明灯效果很差，
        却是很好的小型加热器。
  \item 钨丝温度为多少时，$\lambda_{\max}$ 会位于
        $\qty{550}{nm}$（可见光中部）？钨的熔点为
        $\qty{3700}{K}$；总结为什么白炽灯无法模拟
        日光，以及照明为何转向其他技术。
\end{enumerate}
\end{exercise}

\begin{exercise}[$\star\star\star$]\label{exo:g10:light-spectra:13}
恒星的光谱为连续光谱，峰值位于
$\qty{420}{nm}$，其上有暗线，位于 $393$、$397$、$410$、$434$、
$486$ 和 $\qty{656}{nm}$ 处。对照表：氢
$\{410, 434, 486, 656\}\,\unit{nm}$；钙 $\{393, 397\}\,\unit{nm}$；
钠 $\{589\}\,\unit{nm}$。
\begin{enumerate}
  \item 计算该恒星的表面温度。它比太阳更热
        还是更冷？
  \item 其大气中一定含有哪些元素？
  \item 该大气中钠是否丰富？请说明理由。
\end{enumerate}
\end{exercise}

\begin{exercise}[$\star\star\star$]\label{exo:g10:light-spectra:14}
某放电灯内含两种气体的混合物，其光谱显示
谱线位于 $410$、$434$、$447$、$486$、$502$、$588$、$656$ 和
$\qty{668}{nm}$ 处。利用\cref{exo:g10:light-spectra:8}中的对照表：
\begin{enumerate}
  \item 指出这两种气体。
  \item 一位同学认为，$\qty{588}{nm}$ 谱线证明该灯还含有钠
        （谱线位于 $\qty{589}{nm}$），因为廉价的分光镜无法分辨
        相距 $\qty{1}{nm}$ 的波长。
        在没有更好仪器的情况下，为何氦足以解释这条谱线，
        以及需要怎样的测量，
        才能彻底判定钠是否存在？
\end{enumerate}
\end{exercise}

\begin{exercise}[$\star\star\star$]\label{exo:g10:light-spectra:15}
在猎户座中，参宿四呈明显的红色，参宿七呈
蓝白色。两者连续光谱的峰值分别约位于 $\qty{850}{nm}$ 和
$\qty{240}{nm}$。
\begin{enumerate}
  \item 计算两颗恒星的表面温度。
  \item 不重新计算任何一个温度，直接由两个
        波长求出两者的温度之比。
  \item 已知\emph{两个峰值都不}位于
        可见光范围内，试解释这两种颜色。
  \item 两颗恒星肉眼看来都很亮。为什么峰值位于
        可见光范围之外，却不会使恒星不可见？
\end{enumerate}
\end{exercise}

\section{问题：读取星光}

\begin{problem}[{周末问题 --- 读取星光：说“永远”的哲学家，
以及三棱镜如何测定恒星的温度，并读取
其化学成分}]
\label{pb:g10:light-spectra:1}
1835 年，哲学家奥古斯特·孔德需要一个例子，说明
人类永远无法触及的知识，他的选择颇为审慎：恒星的化学
成分。没有人会把一颗恒星装进瓶子。1859 年，
基尔霍夫与本生把太阳的暗线对应到实验室火焰的亮线上，
这种原本遥不可及的知识，便借由返回的光而到达。
本题将完整重走这一路径：实验室中的光谱
指纹、发光体、太阳的暗线——最后是一颗恒星，
其温度和成分，你将
在书桌前测定。

\medskip
\noindent\textbf{第一部分 --- 存档的指纹。}
\begin{enumerate}
  \item 以纳米为单位给出可见光的波长范围，以及
        两端的颜色。$\qty{550}{nm}$ 位于何处？
  \item 氢放电灯的谱线位于 $410$、$434$、$486$
        和 $\qty{656}{nm}$。给出每条谱线的颜色。
  \item 钠灯基本只发出一条位于 $\qty{589}{nm}$ 的谱线。
        这盏灯在人眼看来呈什么样子？为何在这种路灯下
        判断停放汽车的颜色注定徒劳？
  \item 解释为什么一组谱线能够鉴定一种元素——
        以及为何在判定某元素存在之前，必须找到该元素的
        \emph{全部}强谱线，才能断定其存在。
  \item 预测所见的光谱：(a) 白光穿过一烧瓶
        低温钠蒸气时；(b) 同一蒸气单独置于暗室中，
        受放电激发时所见的光谱。
\end{enumerate}

\medskip
\noindent\textbf{第二部分 --- 热体及其颜色。}
\begin{enumerate}[resume]
  \item 叙述提高一个致密发光体的温度对其连续光谱
        产生的两种影响。
  \item 某卤素灯丝工作于 $\qty{3400}{K}$。计算
        $\lambda_{\max}$，并给出其所属的波段。
  \item 太阳的连续光谱峰值位于 $\qty{500}{nm}$。计算
        其表面温度。
  \item 该峰值位于绿色区域——然而太阳看上去
        并非绿色。请解释原因。
  \item 留在锻炉中的火钳，人在远处先感受到热，却看不到
        它发光；随后发暗红色光，再变为橙白色。
        请结合第 6 题解释整个过程。
\end{enumerate}

\medskip
\noindent\textbf{第三部分 --- 太阳的暗线。}
\begin{enumerate}[resume]
  \item 太阳光谱是一条连续的带，其上穿插着数以千计的
        细小暗线。太阳的哪一部分产生连续背景，
        哪一部分产生谱线？
  \item 太阳光谱中的强谱线位于 $393$、$397$、$410$、$434$、$486$、
        $589$ 和 $\qty{656}{nm}$。对照表：氢
        $\{410, 434, 486, 656\}\,\unit{nm}$；钠 $\{589\}\,\unit{nm}$；
        钙 $\{393, 397\}\,\unit{nm}$。太阳大气中
        存在哪些元素？
  \item 氢的谱线远比其他谱线强。这说明
        太阳的组成有何特点，
        现代测量又给出了怎样的结论？
  \item 1868 年，在太阳大气中观测到一条位于 $\qty{587.6}{nm}$ 的谱线，
        它与任何对照表都不相符。请重构当时的推理，
        以说明为何足以宣布发现新元素，并讲述此事的结局。
  \item 为何太阳的谱线呈暗色，而同样的原子在
        实验室放电中却给出亮线？在日全食期间，
        仅有数秒钟，只有太阳稀薄的大气层仍可在暗色天空
        的衬托下被看见——同样的谱线闪烁为
        \emph{亮}线。请解释原因。
\end{enumerate}

\medskip
\noindent\textbf{第四部分 --- 书桌上的一颗恒星。}
摆在你面前的是亮星织女星的光谱：一条连续光谱，
峰值位于 $\qty{290}{nm}$，其上有强暗线，位于
$410$、$434$、$486$ 和 $\qty{656}{nm}$ 处。
\begin{enumerate}[resume]
  \item 计算织女星的表面温度。
  \item 峰值位于紫外区。预测该恒星看上去的
        颜色，并协调“峰值位于可见光范围之外”与
        “亮星”之间的矛盾。
  \item 指出在织女星大气中占主导地位的元素。
  \item 比较织女星与太阳：温度之比、颜色，以及
        单位发光表面的亮度。
  \item 终章——填写孔德曾断言不可能获得的
        “身份表”：温度、主导元素与织女星看上去的颜色，
        全部由同一束光读出。“永远”这一说法
        持续了多少年？
\end{enumerate}
\end{problem}
